\scanpage{210}
Si p. ex. $R$ est le graphe d'une \uline{bijection} entre $S$ et $S^\vee$, alors \deleted{\unsure{on|on a}} $Z_S=\unsure{\mathfrak S|\mathscr S}_S$, $Z_{S^\vee}=\mathfrak S_{S^\vee}$, et identifiant $S^\vee$ à $S$ par $R$, \inserted{d'où} $\mathcal H\subset\mathfrak S_S\times\mathfrak S_S$, $\mathcal H$ s'identifie à l'antidiagonale

\[
\mathcal H=\{(u,u^{-1})\in\mathfrak S_S\times\mathfrak S_S\mid u\in\mathfrak S_S\}, \tag*{(}
\]

qui n'est pas un \unsure{ss-groupe|ss-gpe|sous-groupe} \deleted{si $\operatorname{card}S\ge3$} \unsure{ssi.|dès que} $\mathfrak S_S$ n'est pas comm i.e. $\operatorname{card}S\ge3$. La définition \inserted{de $\mathcal H$} peut paraître saugrenue, pourtant elle est fortement suggérée par les considérations précédentes.

Ce n'est pas simplement un canular de notation -- ainsi, rien ne nous empêcherait de modifier la définition en apparence saugrenue des couples $(\lambda,\lambda^\vee)$ qualifiés d'homothéties, en les remplaçant par les couples $(\lambda^{-1},\lambda^\vee)$ ou $(\lambda,\lambda^{\vee-1})$ -- par les automorphismes du graphe, ce qui reviendrait à changer de notation pour les relations entre $\varphi_0,\varphi$ via $\lambda,\lambda^\vee$, en écrivant

\[
\deleted{\unsure{\lambda}}\
\left\{\begin{aligned}
&\varphi=\lambda^\vee\varphi_0=\varphi_0\lambda^{\unsure{\vee-1|-1}}
[=(\lambda\varphi_0^{-1})^{-1}]\\
&\deleted{\text{donc}}\quad\deleted{\varphi_0(\lambda}
\left\{\begin{aligned}&\lambda^\vee=\varphi_0(\lambda)^{-1}\\&\lambda=(\varphi_0^{-1}(\lambda^\vee))^{-1}\end{aligned}\right.\\
&\varphi_0(\lambda^{-1}s)=\lambda^\vee\varphi_0(s)\quad\text{ou }\varphi_0(\lambda s)=\lambda^{\vee-1}\varphi_0(s)
\end{aligned}\right.
\]

% slanted
\margin{Cette définition de $\lambda,\lambda^\vee$ en termes de $\varphi_0,\varphi$ a aussi l'avantage d'être stable par symétrie, qui échange alors $\lambda,\lambda^\vee$.}

ce qui donne par contre

\[
(\lambda,\lambda^\vee)(R)=R\quad\text{i.e.}\quad
\inserted{\forall s,s^\vee\in S\times S^\vee}\ (s,s^\vee)\in R\iff(\lambda s,\lambda^\vee s^\vee)\in R.
\]

Dans le cas de l'\unsure{exemple|example} linéaire, ça signifie que si on fait opérer les scalaires \inserted{$\in\Lambda^*$} sur $M$ de la façon habituelle, on les fait opérer sur $M^{\unsure{\vee|*}}$ par contragrédience (la ``multiplication'' inhabituelle), ce qui a l'avantage

% archive sheet number: 209
